% !PRINTED 356
\weberpage{356}

%<<CONT
\gesperrt{considered, then $\mu$ is called a basis form of the ideal $\mathfrak{m}$ in the order $[Q]$.}

For the proof we observe that, by (4), every number of the order $[Q]$ is contained in the form
\weq{10}{\omega = z + \frac{b + \sqrt{D}}{2} y,}
where $y$ and $z$ are rational integers. Now if $\mathfrak{m}$ is primary, then
\weq{11}{N(\mathfrak{m}) = a}
is the least natural number divisible by $\mathfrak{m}$, and hence if $\omega$ is to be divisible by $\mathfrak{m}$, then $z$ must be divisible by $a$. If therefore we set $z = ax$, we obtain the form (9). Conversely it is evident that every number of this form is divisible by $\mathfrak{m}$.

That the representation of a number $\mu$ in this form is possible in only one way follows from the fact that $\sqrt{D}$ is irrational.

The form $\mu$ is therefore called a \gesperrt{basis form of the ideal $\mathfrak{m}$ in the order $Q$}. From it one obtains a basis form of the same ideal in the field $\Omega$:
\[ \lambda = ax + \frac{b' + \sqrt{\Delta}}{2} y, \]
if $b'$ is determined from the congruence
\[ b'Q \equiv b \pmod{2a} \]
which is always solvable, since $Q$ is relatively prime to $a$, and for even $Q$ also $b$ is even.

From (9) there results
\weq{12}{N(\mu) = a(ax^2 + bxy + cy^2).}
The number $\mu$ which one obtains when definite numbers are put in (9) for $x, y$ we can decompose into ideal factors:
\[ \mu = \mathfrak{a}\mathfrak{m}, \]
if
\weq{13}{N(\mathfrak{a}) = ax^2 + bxy + cy^2}
holds, and the following theorem is valid:

8. Theorem. \gesperrt{If $\mu$ is relatively prime to $Q$, and $x, y$ have no common divisors, then the ideal $\mathfrak{a}$ is primary.}

We prove it thus: Let $p$ be a rational prime dividing $\mathfrak{a}$; then $\mu$ is divisible by $p$, and from this it follows first of all that $y$ must be divisible by $p$; for a number $\omega$ of the form (10) can be divisible by a prime $p$ that does not divide $Q$ only if both $y$ and $z = ax$ are divisible by $p$. Accordingly, since by assumption $x$ and $y$
%CONT>>
